% Méthode : cible, variables, protocole, modèles (parties 3 et 4 ; ADR-0009 à 0012, 0031 à 0039).
% Équations détaillées : docs/latex/mathematiques/.

\chapter{Méthode}
\label{chap:methode}

\section{Cible, sortie et horizon}

La cible est le nombre de buts de \emph{chaque équipe} ; la sortie publiée est la loi du \emph{total}
\(T\), repliée sur \(\{0, \dots, 9, \ge 10\}\) (ADR-0009). Modéliser par équipe donne la loi du total
par convolution, et aussi le score exact et le 1N2. L'horizon du MVP, \og H1 \fg{}, est la veille du
match, avant la publication des compositions : seules des informations connues à ce moment entrent dans
les variables (ADR-0010). L'horizon \og H2 \fg{}, une heure avant le coup d'envoi (\og T-60 \fg{}), connaît
en plus le onze de départ et l'entraîneur ; il fait l'objet de la version intermédiaire et se place en
rejeu seulement, faute de compositions après la fin de l'abonnement.

\section{Variables}

Les variables sont regroupées par information, pour mesurer l'apport de chaque groupe :

\begin{center}
\begin{tabularx}{\textwidth}{lX}
\toprule
Groupe & Contenu \\
\midrule
G0, contexte & domicile, championnat, huis clos (périodes sourcées, ADR-0033) \\
G1, force & Elo maison : une échelle par pays, mise à jour par jour, paramètres réglés sur 2005-06 à
            2014-15 puis figés (ADR-0031) \\
G2, production & moyennes glissantes des buts et de l'\code{xg\_proxy}, marqués et encaissés, de
                 l'équipe et de l'adversaire, à décroissance exponentielle en jours et tirées vers la
                 moyenne du championnat (ADR-0032) \\
G3, calendrier & jours de repos, matchs des 14 derniers jours, match européen dans les 4 jours \\
G4, entraîneur (H2) & ancienneté de l'entraîneur dans le club, plafonnée à 365 jours ; indicateur
                      \og nouvel entraîneur \fg{} (moins de 5 matchs) \\
G5, qualité du onze (H2) & note moyenne des 11 titulaires, chaque note étant tirée vers la moyenne de son
                           poste selon le nombre de minutes jouées ; part des \og habituels \fg{}
                           absents ; âge moyen \\
G6, stabilité du onze (H2) & nombre de titulaires communs avec le match précédent ; cohésion par paires
                             pondérée par le rang des matchs ; nouveaux titulaires \\
\bottomrule
\end{tabularx}
\end{center}

L'\code{xg\_proxy} remplace l'xG, absent d'API-FOOTBALL avant 2022-23 : une combinaison des tirs cadrés
et non cadrés, à coefficients positifs ou nuls, estimée une seule fois par moindres carrés sur 2015-16 à
2020-21 (avant tous les plis) puis figée. Le coefficient des tirs non cadrés sort nul :
\(\code{xg\_proxy} = 0{,}3076 \times\) tirs cadrés (ADR-0032, ADR-0035). Sur 2022-23 à 2024-25, sa
corrélation avec l'xG d'API-FOOTBALL valait 0,685 par équipe et par match (première estimation,
ADR-0032). Les moyennes
glissantes ne repartent pas de zéro à chaque saison : l'historique compte sur 730 jours, avec une
demi-vie choisie par validation interne (60, 120 ou 240 jours).

\section{Protocole d'évaluation}

Le protocole est écrit et figé avant toute évaluation (ADR-0037, tag \code{protocole-v1}).

\begin{itemize}
  \item \textbf{Plis} : quatre plis à fenêtre croissante. Pour chaque saison de test, de 2021-22 à
        2024-25, apprentissage sur toutes les saisons antérieures depuis 2015-16, évaluation sur les
        matchs de championnat du top 5 : 7\,156 matchs au total. Les réglages se choisissent par
        validation interne, sur la dernière saison d'apprentissage.
  \item \textbf{Métrique principale} : log-loss de la loi du total ; en complément RPS, Brier de
        \(P(T > 2{,}5)\), calibration par décile, couverture de l'intervalle \([q_{10} ; q_{90}]\).
  \item \textbf{Comparaisons} : sur les mêmes matchs, par différence de log-loss ; intervalle à 95~\%
        par bootstrap par blocs (championnat, saison, journée), 10\,000 rééchantillonnages, graine
        fixe ; test de Diebold-Mariano.
  \item \textbf{Règle de décision} : un modèle plus complexe remplace le plus simple si (i) son gain
        poolé est positif avec un intervalle qui exclut 0, (ii) il gagne dans au moins 3 plis sur 4,
        (iii) sa calibration ne se dégrade pas. Sinon, on garde le plus simple.
  \item \textbf{Essais comptés} : chaque expérience, échecs compris, laisse un rapport ; l'index en
        donne le nombre (17 essais à ce jour, 0 échec).
\end{itemize}

L'intervalle \([q_{10} ; q_{90}]\) d'une loi discrète ne couvre pas exactement 80~\% des matchs : chaque
prédiction publie donc sa \emph{couverture annoncée}, et le critère compare la couverture observée à la
moyenne des couvertures annoncées, à \(\pm 3\) points (ADR-0009).

\section{Modèles}

Les modèles suivent une progression où chaque étape répond à une question (rapport de cadrage, I.4) :

\begin{center}
\begin{tabularx}{\textwidth}{lX}
\toprule
Modèle & Question \\
\midrule
B0 & référence : même loi de Poisson pour tous les matchs \\
B1 & référence : moyennes de buts à domicile et à l'extérieur du championnat \\
M1 & une régression linéaire sur le total suffit-elle ? \\
M2 & une loi de comptage (Poisson sur le total) fait-elle mieux ? \\
M3 & faut-il décomposer par équipe (Poisson par équipe) ? \\
M4 & la loi de Poisson est-elle adéquate, ou faut-il une binomiale négative ? \\
M5 & la dépendance entre les deux équipes (Dixon-Coles) change-t-elle la loi du total ? \\
M6 & plus de variables sous régularisation (ridge, élastique net) ? \\
M7 & un boosting d'arbres à perte de Poisson (H2) fait-il mieux que M3 ? \\
\bottomrule
\end{tabularx}
\end{center}

Le marché des paris (probabilité implicite de plus de 2,5 buts, marge retirée) sert de troisième
référence, sur cet événement seulement ; la cote de clôture, qui intègre les compositions, n'est
qu'une borne haute, jamais une référence de l'horizon H1 (ADR-0036).

\section{Test scellé}

La liste du test est figée avant son lancement, avec le code (tag \code{pre-scelle-h1}, ADR-0039) : le
modèle retenu, B0, B1 et le marché, appris sur 2015-16 à 2024-25, évalués une seule fois sur les matchs
scellés. La commande refuse elle-même un second test et tout code modifié depuis le tag.

\section{L'horizon H2 : méthode}
\label{sec:methode-h2}

La version intermédiaire répond à une question : que rapporte la composition des équipes ? La méthode est
celle de l'horizon H1, avec quatre précautions propres.

\begin{itemize}
  \item \textbf{Spécification avant mesure.} Les variables, les grilles, le chemin de décision et ce qui est
        publié sont écrits dans une ADR \emph{avant} chaque expérience (ADR-0048, ADR-0049). La règle de
        décision du protocole figé s'applique mécaniquement, dans l'ordre stabilité, qualité du onze,
        entraîneur, chaque groupe s'ajoutant au modèle retenu juste avant lui.
  \item \textbf{Même fonction, composition à T-60.} Les variables du match cible n'utilisent que son onze
        et son entraîneur annoncés ; les minutes et les notes de ce match n'existent pas encore. Celles des
        matchs \emph{passés} sont prises à un jour strictement antérieur (ADR-0047). Une valeur incalculable
        reste vide, et la ligne est exclue et comptée.
  \item \textbf{Rodage.} Les notes de 2015-16 n'ont pas d'historique pour être tirées vers la moyenne du
        poste : l'apprentissage de la qualité du onze commence en 2016-17, celui de l'entraîneur en 2017-18,
        et ces lignes de rodage sont comptées, non utilisées.
  \item \textbf{Garde de l'horizon H1.} Aucun modèle H1 ni aucune colonne H1 ne change : une PR qui touche
        aux variables ou aux modèles montre que les prédictions du modèle final H1 et les colonnes H1 du jeu
        sont inchangées (ADR-0044).
\end{itemize}

Les trois groupes s'ajoutent à la structure M3 (Poisson par équipe), comme pour H1. Deux autres structures
sont ensuite comparées à M3 avec les variables retenues : un ridge (M6) et un boosting d'arbres à perte de
Poisson (M7), chacune une fois, sur un chemin de décision écrit d'avance (ADR-0049).

\paragraph{Test scellé de H2.} Le second test réutilise les mêmes matchs scellés, une dernière fois, sur une
liste figée par un tag (\code{pre-scelle-h2}, ADR-0050) : le modèle H1 (référence), le modèle H2 retenu, B1
et le marché. La comparaison de décision est unique : modèle H1 contre modèle H2. Avant l'évaluation, les
colonnes H1 du jeu scellé de H2 doivent être identiques au bit près à celles du premier test ; le nombre de
lignes écartées faute de variables est publié ; un échec technique ne se relance pas sans nouvelle ADR. Les
règles de lecture sont écrites avant le test : un écart non significatif, résultat attendu vu la puissance, ne
retire pas le groupe.

\paragraph{Après le test.} Le rejeu des saisons scellées s'ouvre (ADR-0052) sur trois règles : la condition
est lue dans le journal du scellé, le modèle est celui du test, le score est servi mais pas les cotes. Une
commande de garde vérifie que le rejeu reproduit les prédictions des deux tests à \(10^{-12}\) près. Le
journal des prédictions live fait l'objet d'un relevé \emph{descriptif} : aucune règle de décision, aucun
modèle changé à partir de lui, aucune conclusion avant 8 journées.
